% Supplementary Material for the ArtiGate-YOLO manuscript.
% Section, table, figure and equation numbers carry the prefix S. Every value written as \res or
% \TBD{...} does not exist yet; fill it from runs/ (see paper/README.md for the script behind each table).
% Values that are not marked are design constants read from the code on branch feat/artigate-v0.1.

\documentclass[10pt,a4paper]{article}
\usepackage[margin=2cm]{geometry}
\usepackage{amsmath,amssymb}
\usepackage{graphicx}
\usepackage{booktabs}
\usepackage{multirow}
\usepackage{array}
\usepackage{tabularx}
\usepackage{longtable}
\usepackage[labelfont=bf,font=small]{caption}
\usepackage{float}    % [H]: every table and figure stays in its section
\usepackage{tocloft}
\setlength{\cftsecnumwidth}{2.8em}
\setlength{\cftsubsecnumwidth}{3.6em}
\usepackage[hidelinks]{hyperref}
% Shared macros for the main manuscript and the Supplementary Material.
\usepackage{xcolor}
% \TBD{...} marks every value that does not exist yet (results, decisions, citations to check).
% Search the .tex files for "\TBD" before submission; the PDF shows them in red.
\newcommand{\TBD}[1]{\textcolor{red}{[#1]}}
\newcommand{\res}{\TBD{--}} % one result cell that the experiments will fill
\newcommand{\etal}{\textit{et~al.}}
\newcommand{\method}{ArtiGate-YOLO}
\newcommand{\aga}{AGA}
\newcommand{\Abar}{\bar{A}}
% Placeholder for a figure that will be produced from the runs.
\newcommand{\placeholderfig}[2]{%
  \fbox{\parbox[c][#1][c]{0.94\linewidth}{\centering\small\textcolor{red}{#2}}}}


\renewcommand{\thesection}{S\arabic{section}}
\renewcommand{\thetable}{S\arabic{table}}
\renewcommand{\thefigure}{S\arabic{figure}}
\renewcommand{\theequation}{S\arabic{equation}}

% Standard note under every table with paired statistics.
\newcommand{\pairedcaveat}[1]{%
  \par\smallskip{\footnotesize\noindent Differences are row minus #1, per image after averaging each image over the three seeds; 95\% percentile bootstrap interval over images (2,000 resamples, seed 0); $p$: two-sided Wilcoxon signed-rank test with zero differences discarded; $p_{\mathrm{Holm}}$: Holm adjustment within the family named in the caption. W/T/L: images on which the row is better / equal / worse than #1.\par}}

\title{Supplementary Material for\\[2pt]
\large ArtiGate-YOLO: Artifact-Supervised Gated Attention for Robust Real-Time Skin Lesion Segmentation}
\author{Nguyen Chi Hai, Le Huu Son Hai, Nguyen Minh Phuc and Ngo Hoang Tu\\[2pt]
\small Faculty of Information Technology, Van Lang School of Technology, Van Lang University, Ho Chi Minh City 70000, Vietnam\\
\small Corresponding author: Ngo Hoang Tu (tu.nh@vlu.edu.vn)}
\date{}

\begin{document}
\maketitle

\noindent This document accompanies the main manuscript. It gives the metric conventions (S1), training histories (S2), the parameters of the artifact generator and the partial labels (S3--S4), the complete statistical tables behind the main results (S5--S9, S11), cost measurements (S10), published values of related methods (S12), implementation details (S13), the data inventory (S14) and the relation to our conference paper (S15). Unless stated otherwise, results are on the official ISIC 2018 Task~1 test set (1,000 images) on the $640\times640$ letterbox canvas with the top-1 mask rule.

\tableofcontents
\clearpage

% =================================================================================================
\section{Metric Definitions and Mask Rules}\label{sec:s-metrics}
All metrics are computed per image on the $640\times640$ canvas, on which images and ground-truth masks are letterboxed with the same transform. None of the four metrics counts true negatives, so the padded area does not inflate them. Let $P$ and $G$ be the predicted and ground-truth pixel sets.
\begin{align}
\mathrm{Dice}&=\frac{2|P\cap G|}{|P|+|G|}, &
\mathrm{IoU}&=\frac{|P\cap G|}{|P\cup G|}.
\end{align}
The boundary $\partial M$ of a mask is the set of mask pixels removed by one binary erosion with the 4-connected cross (pixels outside the image count as background). With $d(x,\partial M)$ the Euclidean distance from pixel $x$ to the nearest pixel of $\partial M$ (exact distance transform),
\begin{equation}
\mathrm{HD95}=Q_{0.95}\big(\{d(x,\partial G):x\in\partial P\}\cup\{d(y,\partial P):y\in\partial G\}\big),
\end{equation}
the 95th percentile of the pooled symmetric boundary distances (MedPy convention), in canvas pixels. The boundary F1 score~\cite{csurka2013bf} uses a tolerance $\tau=0.0075\times$ the image diagonal, $\tau=6.79$~px on the canvas:
\begin{equation}
\mathrm{BP}=\frac{|\{x\in\partial P:d(x,\partial G)\le\tau\}|}{|\partial P|},\quad
\mathrm{BR}=\frac{|\{y\in\partial G:d(y,\partial P)\le\tau\}|}{|\partial G|},\quad
\mathrm{BF1}=\frac{2\,\mathrm{BP}\,\mathrm{BR}}{\mathrm{BP}+\mathrm{BR}}.
\end{equation}

\paragraph{Empty masks}
If both masks are empty, Dice, IoU and BF1 are 1 and HD95 is 0. If exactly one is empty, Dice, IoU and BF1 are 0 and HD95 is the canvas diagonal (905.1~px), the largest possible distance. Every ISIC 2018 test image contains a lesion, so an empty prediction always receives these worst-case values; their number per model is reported in Table~\ref{tab:s-empty}.

\paragraph{From instances to one mask}
The models predict lesion instances. Instance masks are produced at input resolution (\texttt{retina\_masks}) and binarized at probability 0.5. The \emph{top-1} rule keeps the mask of the instance with the highest confidence among all predictions with confidence $\ge0.001$; it is the rule of the conference paper and the primary analysis. The \emph{union} rule merges every instance with confidence $\ge0.25$; it is a sensitivity analysis (Section~\ref{sec:s-sensitivity}). Both rules and the non-inferiority margin ($\delta=0.5$ Dice points) were fixed on 25 October 2026, before any ISIC 2018 test image was scored.

\paragraph{Detection metrics}
Mask mAP50 and mAP50--95 are computed by the Ultralytics 8.4.13 validator on the same split and are point estimates without intervals.

\begin{table}[H]
\centering
\caption{Images with an empty predicted mask (worst-case scores) per model and seed, ISIC 2018 test set, top-1 rule.}
\label{tab:s-empty}
\small
\begin{tabular}{lccc}
\toprule
Model & Seed 0 & Seed 1 & Seed 2 \\
\midrule
A1 & \res & \res & \res \\
A1s & \res & \res & \res \\
A4 & \res & \res & \res \\
\TBD{all other rows} & \res & \res & \res \\
\bottomrule
\end{tabular}
\end{table}

% =================================================================================================
\section{Training Histories}\label{sec:s-histories}
\begin{figure}[H]
\centering
\placeholderfig{6cm}{Training curves of A1, A1s, A3 and A4 (seed 0): total segmentation loss, artifact loss (A2--A7), validation mask mAP50--95 and the learned $\beta_2,\beta_3$ per epoch. Source: runs/segment/<run>/results.csv and the AGA parameters of the saved checkpoints.}
\caption{Training histories. \TBD{Epoch of the best checkpoint and of early stopping for every run (Table~\ref{tab:s-seeds}).}}
\label{fig:s-histories}
\end{figure}

% =================================================================================================
\section{Synthetic Artifact Generator}\label{sec:s-generator}
The generator (\texttt{src/artigate/data/synth\_artifacts.py}) draws on 8-bit BGR images and returns, for every type, the \emph{core mask} used as the present label and the set of all pixels it changed; pixels that were changed but are outside the core (anti-aliased edges, soft fall-off) are set to unknown. All sizes are relative to $s$, the shorter side of the image content area (letterbox padding excluded). Types are drawn in physical order: marker ink, hair, ruler, dark corner.

\begin{table}[H]
\centering
\caption{Generator parameters. $U(a,b)$: uniform; severity $k\in\{1,2,3,4\}$ is drawn per type and image during training.}
\label{tab:s-generator}
\footnotesize
\begin{tabularx}{\textwidth}{@{}l>{\raggedright\arraybackslash}X>{\raggedright\arraybackslash}p{0.27\textwidth}@{}}
\toprule
Type & Geometry and appearance & Severity dependence and label \\
\midrule
Hair &
Cubic B\'ezier strands of length $U(0.15,0.8)\,s$, random direction, start point up to 10\% outside the image, inner control points displaced perpendicular to the chord with a standard deviation of $0.15\times$ the length; width $\max(1,U(0.002,0.007)\,s)$ px; anti-aliased, Gaussian blur $\sigma=0.6$, opacity $U(0.75,1)$. Colour: 90\% dark brown/black (base intensity $U(10,60)$ with per-channel factors), 10\% light/grey ($U(120,190)$). May cross the lesion. &
Strands per image: $U\{2..5\}$, $U\{6..14\}$, $U\{15..34\}$, $U\{35..69\}$ for $k=1..4$. Present: opacity $>0.35$. \\
\addlinespace
Ruler &
Band along one image border at an offset of $U(0,0.04)\,s$, width $\max(6,U(0.04,0.08)\,s\cdot w_k)$ px, covering $U(0.5,1)$ of the border ($U(0.75,1)$ for $k\ge3$). 85\% light body ($U(215,255)$) with dark ticks, 15\% dark body with light ticks; ticks every $U(0.012,0.025)\,s$ (long every 10th, medium every 5th), numerals at 70\% of the long ticks on bands at least 18~px high; opacity $U(0.55,0.9)$. Rejected if it would overlap the lesion box enlarged by $0.03\,s$. &
$w_k=0.6,0.85,1.1,1.25$; two rulers on different borders at $k=4$. Present: the whole band. \\
\addlinespace
Marker ink &
Around the lesion box on an ellipse with semi-axes $1.2$--$1.5\times$ the box half-sizes plus $0.04\,s$. Style: dots (45\%), arc (35\%) or strokes (20\%); line width $\max(2,U(0.008,0.02)\,s)$; colour purple (75\%), blue (15\%) or black (10\%); Gaussian blur $\sigma=U(0.8,2)$, opacity $U(0.5,0.9)$. Nothing is drawn inside the lesion box enlarged by $0.01\,s$, before or after blurring. &
Dots: $U\{2+2k..3+3k\}$; arc span $60^\circ,120^\circ,220^\circ,340^\circ$ ($\pm20^\circ$); strokes: $U\{k..2k+1\}$, length $U(0.04,0.12)\,s$. Present: opacity $>0.2$. \\
\addlinespace
Dark corner &
Circular aperture centred near the image centre (normal jitter of 3\% per axis) with radius $r$ times the half-diagonal and a smoothstep edge of width $U(0.01,0.05)\,s$; outside, the image is replaced by dark noise (level $U(0,25)$, $\sigma=3$). The radius is enlarged if needed so the lesion box stays inside with a margin of $0.03\,s$; if the aperture would then cover the whole image, no corner is drawn. &
$r\sim U(0.92,0.99)$, $U(0.82,0.92)$, $U(0.72,0.82)$, $U(0.62,0.72)$ for $k=1..4$. Present: blend weight $>0.5$. \\
\midrule
Training &
\multicolumn{2}{>{\raggedright\arraybackslash}p{0.68\textwidth}}{Probability that an image receives artifacts $p_{\mathrm{synth}}=0.7$; type probabilities 0.5 (hair), 0.3 (ruler, ink, corner), at least one type per image; severity $U\{1..4\}$ per type; artifacts are drawn inside the letterbox content area only; the lesion box is the union of the ground-truth instance boxes.} \\
Severity sweep &
\multicolumn{2}{>{\raggedright\arraybackslash}p{0.68\textwidth}}{One type forced, or ``mix'' with type probabilities 0.6/0.4/0.4/0.4, at a fixed severity; per-image seed from the MD5 hash of image name, type and severity plus a base of 10,000; JPEG quality 95.} \\
\bottomrule
\end{tabularx}
\end{table}

\noindent Figure~\ref{fig:s-generator} shows \TBD{n} random training samples after the full augmentation pipeline with the four target maps (\texttt{scripts/viz\_train\_batch.py}).
\begin{figure}[H]
\centering
\placeholderfig{5cm}{Augmented training batch (mosaic, rotation, synthetic artifacts) with the four target channels: present (white), absent (black), unknown (grey). Produce with scripts/viz\_train\_batch.py on the ISIC 2018 training set.}
\caption{Training samples and their partial artifact targets after augmentation.}
\label{fig:s-generator}
\end{figure}

% =================================================================================================
\section{Partial Labels and Target Coding}\label{sec:s-labels}
For every training image and artifact type, each pixel is present, absent or unknown. The target maps start as unknown and are filled in the order of Table~\ref{tab:s-labels}; later steps overwrite earlier ones.

\begin{table}[H]
\centering
\caption{Sources of the artifact targets, in the order they are applied.}
\label{tab:s-labels}
\footnotesize
\begin{tabularx}{\textwidth}{@{}l>{\raggedright\arraybackslash}X@{}}
\toprule
Step & Rule \\
\midrule
1. Rules (absent) &
Hair: no thin dark structure, i.e.\ morphological black-hat response of the grey image $\le8$ (elliptical element of $\max(9,0.025\,s')$ px, made odd, with $s'$ the shorter side of the loaded image). Ruler: neither white/grey (HSV saturation $<40$ and value $>150$) nor very dark (value $<60$). Marker ink: neither purple/blue (OpenCV hue 95--170 with saturation $>40$) nor very dark. The candidate regions of these three rules are dilated by a $5\times5$ square before taking the complement. Dark corner: value $\ge70$. Ruler and marker ink are also absent inside the lesion box shrunk by 10\% per side. \\
\addlinespace
2. Image-level labels (absent) &
For ISIC 2018 training images whose annotation~\cite{bissoto2020debiasing} says a type is not present, every pixel of that type is absent. Mapping of the annotation columns: \texttt{hair}$\to$hair, \texttt{ruler}$\to$ruler, \texttt{ink}$\to$marker ink, \texttt{dark\_corner}$\to$dark corner (\texttt{scripts/make\_artifact\_flags.py}). A positive image-level label sets nothing, because it does not say where the artifact is. \\
\addlinespace
3. Synthetic artifacts &
Core mask present; other changed pixels unknown (Section~\ref{sec:s-generator}). \\
\bottomrule
\end{tabularx}
\end{table}

\paragraph{Coding and augmentation}
The four maps are stored as image channels 4--7 with the codes 255 (present), 0 (absent) and 114 (unknown, the Ultralytics padding value) and are transformed with the image. After augmentation a value $v$ is read as present if $v/255\ge0.75$, absent if $v/255\le0.1$ and unknown otherwise. The absent threshold lies well below the unknown code (0.447), so that a pixel mixed by mix-up from an absent and an unknown pixel (about 57, or 0.22) stays unknown. Padding introduced by the augmentations is unknown (114) or absent (0); both are correct because padding contains no artifact. A synthetic artifact drawn later can cover an earlier one (for example, the aperture over hair); the earlier core mask is kept, so a few covered pixels remain labelled present. The colour augmentation (HSV gains) acts on channels 1--3 only.

\paragraph{Downsampling to stride 4}
For the loss, a cell of the $160\times160$ logit grid is positive if any pixel inside it is present (max pooling), negative if every pixel inside it is absent (min pooling), and ignored otherwise. A type contributes to the loss of a batch only if at least one of its cells is labelled.

\paragraph{Coverage}
\TBD{Table~\ref{tab:s-coverage}: share of present, absent and unknown cells per type over one training epoch (seed 0), to show that negatives dominate but do not cover every pixel.}
\begin{table}[H]
\centering
\caption{Label coverage of the artifact targets at stride 4 over one training epoch (\%).}
\label{tab:s-coverage}
\small
\begin{tabular}{lccc}
\toprule
Type & Present & Absent & Unknown \\
\midrule
Hair & \res & \res & \res \\
Ruler & \res & \res & \res \\
Marker ink & \res & \res & \res \\
Dark corner & \res & \res & \res \\
\bottomrule
\end{tabular}
\end{table}

% =================================================================================================
\section{Complete Ablation Results}\label{sec:s-ablation}
Tables~\ref{tab:s-ablation-dice}--\ref{tab:s-ablation-shape} extend the ablation of the main text to all metrics. Each row is compared with A4; the family is the seven comparisons of one metric. Script: \texttt{scripts/compare\_family.py} with A4 as reference.

\begin{table}[H]
\centering
\caption{Ablation, Dice and IoU. Family: rows A0--A7 versus A4 on the full test set, per metric.}
\label{tab:s-ablation-dice}
\footnotesize
\setlength{\tabcolsep}{3.5pt}
\begin{tabular}{lcccccccccc}
\toprule
& \multicolumn{5}{c}{Dice} & \multicolumn{5}{c}{IoU} \\
\cmidrule(lr){2-6}\cmidrule(lr){7-11}
Row & Mean & $\Delta$ [95\% CI] & $p$ & $p_{\mathrm{Holm}}$ & W/T/L & Mean & $\Delta$ [95\% CI] & $p$ & $p_{\mathrm{Holm}}$ & W/T/L \\
\midrule
A0 & \res & \res & \res & \res & \res & \res & \res & \res & \res & \res \\
A1 & \res & \res & \res & \res & \res & \res & \res & \res & \res & \res \\
A1s & \res & \res & \res & \res & \res & \res & \res & \res & \res & \res \\
A2 & \res & \res & \res & \res & \res & \res & \res & \res & \res & \res \\
A3 & \res & \res & \res & \res & \res & \res & \res & \res & \res & \res \\
\textbf{A4} & \res & Ref. & -- & -- & -- & \res & Ref. & -- & -- & -- \\
A5 & \res & \res & \res & \res & \res & \res & \res & \res & \res & \res \\
A7 & \res & \res & \res & \res & \res & \res & \res & \res & \res & \res \\
\bottomrule
\end{tabular}
\pairedcaveat{A4}
\end{table}

\begin{table}[H]
\centering
\caption{Ablation, HD95 (canvas pixels, lower is better) and boundary F1. Family as in Table~\ref{tab:s-ablation-dice}.}
\label{tab:s-ablation-shape}
\footnotesize
\setlength{\tabcolsep}{3.5pt}
\begin{tabular}{lcccccccccc}
\toprule
& \multicolumn{5}{c}{HD95 (px)} & \multicolumn{5}{c}{BF1} \\
\cmidrule(lr){2-6}\cmidrule(lr){7-11}
Row & Mean & $\Delta$ [95\% CI] & $p$ & $p_{\mathrm{Holm}}$ & W/T/L & Mean & $\Delta$ [95\% CI] & $p$ & $p_{\mathrm{Holm}}$ & W/T/L \\
\midrule
A0 & \res & \res & \res & \res & \res & \res & \res & \res & \res & \res \\
A1 & \res & \res & \res & \res & \res & \res & \res & \res & \res & \res \\
A1s & \res & \res & \res & \res & \res & \res & \res & \res & \res & \res \\
A2 & \res & \res & \res & \res & \res & \res & \res & \res & \res & \res \\
A3 & \res & \res & \res & \res & \res & \res & \res & \res & \res & \res \\
\textbf{A4} & \res & Ref. & -- & -- & -- & \res & Ref. & -- & -- & -- \\
A5 & \res & \res & \res & \res & \res & \res & \res & \res & \res & \res \\
A7 & \res & \res & \res & \res & \res & \res & \res & \res & \res & \res \\
\bottomrule
\end{tabular}
\pairedcaveat{A4}
\par\smallskip{\footnotesize\noindent For HD95, W means a lower distance than A4.\par}
\end{table}

\begin{table}[H]
\centering
\caption{External baselines versus A4 on the ISIC 2018 test set. Family: the baselines of the main comparison, per metric.}
\label{tab:s-external-family}
\footnotesize
\setlength{\tabcolsep}{3.5pt}
\begin{tabular}{lcccccc}
\toprule
Model & $\Delta$Dice [95\% CI] & $p_{\mathrm{Holm}}$ & $\Delta$HD95 [95\% CI] & $p_{\mathrm{Holm}}$ & $\Delta$BF1 [95\% CI] & $p_{\mathrm{Holm}}$ \\
\midrule
YOLO26n-seg & \res & \res & \res & \res & \res & \res \\
A1, v7 view & \res & \res & \res & \res & \res & \res \\
A1, v1 + v7 & \res & \res & \res & \res & \res & \res \\
U-Net & \res & \res & \res & \res & \res & \res \\
MALUNet & \res & \res & \res & \res & \res & \res \\
EGE-UNet & \res & \res & \res & \res & \res & \res \\
nnU-Net 2D & \res & \res & \res & \res & \res & \res \\
\bottomrule
\end{tabular}
\pairedcaveat{A4}
\end{table}

\paragraph{Learned gate strengths}
\TBD{$\beta_2$ and $\beta_3$ of every AGA row and seed (mean, min, max), and the norm of the $\bar{A}$-channel weights of the $7\times7$ convolution, to show how far each AGA block moved away from CBAM.}

% =================================================================================================
\section{Artifact-Stratified Results and Test-Set Annotation}\label{sec:s-strata}
\subsection{Annotation of the ISIC 2018 test set}
Two of the authors (L.~H.~S.~Hai and N.~M.~Phuc) marked the presence of hair, ruler, marker ink and dark corner, each on their own computer, with \texttt{scripts/label\_artifacts.py}, which shows one image at a time in identifier order, with one key per type and no model output or mask. The images were prepared with \texttt{scripts/make\_annotator\_bundle.py}: each original image scaled so that its longer side is 900 pixels (area interpolation when reducing, bilinear when enlarging) and saved as JPEG at quality 95, the same view the tool shows when it reads the original files. The annotators saw only these images: no mask, no model output and no test-set result, since no model had yet been scored on the test set. An artifact is marked present if it is visible anywhere in the image at that size: \emph{hair}, at least one hair of any colour, on or off the lesion, including hair under gel; \emph{ruler}, any printed scale, graduation, number or calibration strip (not text added by acquisition software); \emph{marker ink}, any skin-marker trace (dots, circles, strokes), including faint ink, but not the print of a ruler; \emph{dark corner}, an arc-shaped dark region from the dermatoscope aperture in at least one corner (not straight black borders from cropping). Gel bubbles, gel borders and patches are not labelled. Both annotators read the definitions and first labelled the same 20 training images (4 flagged with each type in the training labels of~\cite{bissoto2020debiasing} and 4 with none, drawn with seed 20261026), compared their labels and discussed the differences; these practice labels are not used. They then labelled the test images independently and did not confer until the disagreement step: annotator~1 labelled all 1,000 test images, and annotator~2 a random subset of 200 test images fixed before labelling began (\texttt{data/splits/isic2018\_test\_kappa\_subset.csv}, seed 20261026). Cohen's $\kappa$ (Table~\ref{tab:s-kappa}) is computed on these independent labels. The two annotators then resolved the disagreements on the subset by discussion (\texttt{scripts/artifact\_agreement.py}); the final labels are those of annotator~1 with these resolutions. All labels were committed before any model was scored on the test set \TBD{(commit hash)}; the evaluation script refuses to score a test set until the final labels are committed.

\begin{table}[H]
\centering
\caption{Test-set artifact labels: prevalence and inter-annotator agreement on the common subset (Cohen's $\kappa$~\cite{cohen1960kappa} with a 95\% bootstrap interval over images).}
\label{tab:s-kappa}
\small
\begin{tabular}{lccc}
\toprule
Type & Images with artifact ($n$, \%) & Raw agreement (\%) & $\kappa$ [95\% CI] \\
\midrule
Hair & \res & \res & \res \\
Ruler & \res & \res & \res \\
Marker ink & \res & \res & \res \\
Dark corner & \res & \res & \res \\
Any artifact & \res & -- & -- \\
None & \res & -- & -- \\
\bottomrule
\end{tabular}
\end{table}

\subsection{Complete stratified comparison}
\begin{table}[H]
\centering
\caption{A4 versus A1 by artifact stratum, all metrics. Family (primary): the six strata of this table for Dice; the other metrics are adjusted in their own families. Script: \texttt{scripts/compare\_family.py --strata}.}
\label{tab:s-strata}
\footnotesize
\setlength{\tabcolsep}{3pt}
\begin{tabular}{lccccccccc}
\toprule
Stratum & $n$ & Dice A1 & Dice A4 & $\Delta$Dice [95\% CI] & $p_{\mathrm{Holm}}$ & W/T/L & $\Delta$IoU [95\% CI] & $\Delta$HD95 [95\% CI] & $\Delta$BF1 [95\% CI] \\
\midrule
Any artifact & \res & \res & \res & \res & \res & \res & \res & \res & \res \\
With hair & \res & \res & \res & \res & \res & \res & \res & \res & \res \\
With ruler & \res & \res & \res & \res & \res & \res & \res & \res & \res \\
With marker ink & \res & \res & \res & \res & \res & \res & \res & \res & \res \\
With dark corner & \res & \res & \res & \res & \res & \res & \res & \res & \res \\
Artifact-free & \res & \res & \res & \res & \res & \res & \res & \res & \res \\
\bottomrule
\end{tabular}
\pairedcaveat{A1}
\par\smallskip{\footnotesize\noindent Strata overlap: an image with hair and a ruler is in both. Strata with fewer than 10 images are not tested.\par}
\end{table}

\begin{table}[H]
\centering
\caption{Secondary stratified comparisons (Dice): A4 versus A1s (does supervision add to artifact augmentation?) and A4 versus the best external baseline. Holm within each comparison over the six strata.}
\label{tab:s-strata-secondary}
\footnotesize
\setlength{\tabcolsep}{4pt}
\begin{tabular}{lcccc}
\toprule
& \multicolumn{2}{c}{A4 $-$ A1s} & \multicolumn{2}{c}{A4 $-$ \TBD{best baseline}} \\
\cmidrule(lr){2-3}\cmidrule(lr){4-5}
Stratum & $\Delta$Dice [95\% CI] & $p_{\mathrm{Holm}}$ & $\Delta$Dice [95\% CI] & $p_{\mathrm{Holm}}$ \\
\midrule
Any artifact & \res & \res & \res & \res \\
With hair & \res & \res & \res & \res \\
With ruler & \res & \res & \res & \res \\
With marker ink & \res & \res & \res & \res \\
With dark corner & \res & \res & \res & \res \\
Artifact-free & \res & \res & \res & \res \\
\bottomrule
\end{tabular}
\pairedcaveat{the second model}
\end{table}

% =================================================================================================
\section{Synthetic Severity Sweep}\label{sec:s-severity}
Test images receive artifacts of one type (or a mix) at severity 1--4 (Table~\ref{tab:s-generator}); severity 0 is the original test set. The sweep uses the generator seen in training by A1s--A4 and A7, so it is reported descriptively and is not part of the success criteria.

\begin{table}[H]
\centering
\caption{Dice under synthetic artifacts (mean over images and seeds, 95\% bootstrap interval in the source files). Drop: severity 0 minus severity 4.}
\label{tab:s-severity}
\footnotesize
\setlength{\tabcolsep}{3.5pt}
\begin{tabular}{llcccccc}
\toprule
Type & Model & Sev.\ 0 & Sev.\ 1 & Sev.\ 2 & Sev.\ 3 & Sev.\ 4 & Drop \\
\midrule
\multirow{4}{*}{Hair} & A1 & \res & \res & \res & \res & \res & \res \\
 & A1s & \res & \res & \res & \res & \res & \res \\
 & A4 & \res & \res & \res & \res & \res & \res \\
 & \TBD{best baseline} & \res & \res & \res & \res & \res & \res \\
\midrule
\multirow{4}{*}{Ruler} & A1 & \res & \res & \res & \res & \res & \res \\
 & A1s & \res & \res & \res & \res & \res & \res \\
 & A4 & \res & \res & \res & \res & \res & \res \\
 & \TBD{best baseline} & \res & \res & \res & \res & \res & \res \\
\midrule
\multirow{4}{*}{Marker ink} & A1 & \res & \res & \res & \res & \res & \res \\
 & A1s & \res & \res & \res & \res & \res & \res \\
 & A4 & \res & \res & \res & \res & \res & \res \\
 & \TBD{best baseline} & \res & \res & \res & \res & \res & \res \\
\midrule
\multirow{4}{*}{Dark corner} & A1 & \res & \res & \res & \res & \res & \res \\
 & A1s & \res & \res & \res & \res & \res & \res \\
 & A4 & \res & \res & \res & \res & \res & \res \\
 & \TBD{best baseline} & \res & \res & \res & \res & \res & \res \\
\midrule
\multirow{4}{*}{Mix} & A1 & \res & \res & \res & \res & \res & \res \\
 & A1s & \res & \res & \res & \res & \res & \res \\
 & A4 & \res & \res & \res & \res & \res & \res \\
 & \TBD{best baseline} & \res & \res & \res & \res & \res & \res \\
\bottomrule
\end{tabular}
\end{table}

% =================================================================================================
\section{Artifact Head Evaluation}\label{sec:s-arthead}
The head is scored on \TBD{100} ISIC 2018 validation images with held-out synthetic artifacts (severity 1--4, per-image seeds from a base of 20,000) at a probability threshold of 0.5 (\texttt{scripts/eval\_artifact\_head.py}). IoU is the summed intersection over the summed union of all images; Dice is the mean per image; anti-aliased edges are ignored. ``Clean activation'' is the share of pixels predicted as artifact on the unmodified images, averaged over the four types. For A3 the maps are not trained to match artifact types, so its per-type values show only whether the unsupervised map happens to align with them. Real artifacts in the images count as errors in both measures, so the values are lower bounds that serve to compare rows.

\begin{table}[H]
\centering
\caption{Artifact head by type: IoU / per-image Dice on held-out synthetic artifacts, and clean activation (\%). Mean over three seeds.}
\label{tab:s-arthead}
\footnotesize
\setlength{\tabcolsep}{4pt}
\begin{tabular}{lccccc}
\toprule
Row & Hair & Ruler & Marker ink & Dark corner & Clean act.\ (\%) \\
\midrule
A2 & \res~/~\res & \res~/~\res & \res~/~\res & \res~/~\res & \res \\
A3 (unsupervised) & \res~/~\res & \res~/~\res & \res~/~\res & \res~/~\res & \res \\
\textbf{A4} & \res~/~\res & \res~/~\res & \res~/~\res & \res~/~\res & \res \\
A5 ($p_{\mathrm{synth}}=0$) & \res~/~\res & \res~/~\res & \res~/~\res & \res~/~\res & \res \\
A7 & \res~/~\res & \res~/~\res & \res~/~\res & \res~/~\res & \res \\
\bottomrule
\end{tabular}
\end{table}

\begin{table}[H]
\centering
\caption{A4 artifact head by severity: IoU per type.}
\label{tab:s-arthead-severity}
\small
\begin{tabular}{lcccc}
\toprule
Type & Sev.\ 1 & Sev.\ 2 & Sev.\ 3 & Sev.\ 4 \\
\midrule
Hair & \res & \res & \res & \res \\
Ruler & \res & \res & \res & \res \\
Marker ink & \res & \res & \res & \res \\
Dark corner & \res & \res & \res & \res \\
\bottomrule
\end{tabular}
\end{table}

\paragraph{Real artifacts}
Marker and ruler maps are compared with the FEDD masks~\cite{carrion2023fedd} on the 333 DDI images that have them, and the hair map with the 500 real hair masks of~\cite{hossain2023hairmask}. Matching image identities with \texttt{scripts/check\_leakage.py} places 478 of these images in the ISIC 2018 training set, 15 in its test set and 7 in neither. The model is run on the images of the hair-mask set itself, so each mask is aligned with its own image. The 478 training images were seen during training, but no real hair pixel in them was labelled present; real hair could only be labelled absent by the rules of Section~\ref{sec:s-labels} or by an image-level flag of no hair. We report the 22 unseen images and the 478 seen images separately, and the share of real hair pixels in the 478 that were labelled absent during training (\texttt{scripts/audit\_hair\_negatives.py}). The images of the hair-mask set are the ISIC images stretched to $512\times512$ without cropping; mapped back onto the training canvas, 467 of the 478 seen images contain hair and align with their training image (correlation $\ge0.991$). On these, the hair rule labelled 3.3\% of all real hair pixels absent (median 1.2\% per image, 90th percentile 21\%; 340 of the 467 images below 5\%). The images with the largest shares show very thin, faint hair, which the head may therefore learn to ignore, in unseen images as well. Although every one of these images shows hair, 30 of the 467 carry an image-level flag of no hair~\cite{bissoto2020debiasing}, so all their hair pixels were labelled absent. The rule was fixed before this audit and was not changed in response to it. The 15 test images are scored only with the final checkpoints, like the rest of the test set. None of these masks is used for training or model selection. The images are letterboxed to $640\times640$ like all inputs; each predicted probability map is upsampled to the canvas and scored inside the image content area only. FEDD has no hair or dark-corner class, so only the marker and ruler maps are scored there; FEDD's marker class is compared with our marker-ink map.

\begin{table}[H]
\centering
\caption{Artifact head on real artifacts: IoU, Dice (threshold 0.5) and pixel-level average precision (AP). Mean over three seeds.}
\label{tab:s-arthead-real}
\footnotesize
\setlength{\tabcolsep}{3.5pt}
\begin{tabular}{lcccccccccccc}
\toprule
& \multicolumn{3}{c}{Marker (FEDD)} & \multicolumn{3}{c}{Ruler (FEDD)} & \multicolumn{3}{c}{Hair, unseen (22)} & \multicolumn{3}{c}{Hair, seen (478)} \\
\cmidrule(lr){2-4}\cmidrule(lr){5-7}\cmidrule(lr){8-10}\cmidrule(lr){11-13}
Row & IoU & Dice & AP & IoU & Dice & AP & IoU & Dice & AP & IoU & Dice & AP \\
\midrule
A2 & \res & \res & \res & \res & \res & \res & \res & \res & \res & \res & \res & \res \\
A3 & \res & \res & \res & \res & \res & \res & \res & \res & \res & \res & \res & \res \\
\textbf{A4} & \res & \res & \res & \res & \res & \res & \res & \res & \res & \res & \res & \res \\
A7 & \res & \res & \res & \res & \res & \res & \res & \res & \res & \res & \res & \res \\
\midrule
Images with the artifact ($n$) & \multicolumn{3}{c}{\res} & \multicolumn{3}{c}{\res} & \multicolumn{3}{c}{22} & \multicolumn{3}{c}{478} \\
Hair pixels labelled absent in training (\%) & \multicolumn{3}{c}{--} & \multicolumn{3}{c}{--} & \multicolumn{3}{c}{--} & \multicolumn{3}{c}{3.3 by rules$^{\ast}$} \\
\bottomrule
\end{tabular}
\par\smallskip{\footnotesize $^{\ast}$On the 467 seen images that contain hair and align with their training image; image-level flags add 30 images whose hair is entirely labelled absent.\par}
\end{table}

\noindent\TBD{Agreement of the head with the real test-set labels: image-level AUROC of the maximum of each predicted map against the labels of Section~\ref{sec:s-strata} (no training signal of this kind was used).}

% =================================================================================================
\section{Sensitivity Analyses}\label{sec:s-sensitivity}
\begin{table}[H]
\centering
\caption{Robustness of the primary comparison (A4 $-$ A1, Dice, any-artifact and artifact-free strata) to analysis choices.}
\label{tab:s-sensitivity}
\footnotesize
\setlength{\tabcolsep}{4pt}
\begin{tabular}{lcccc}
\toprule
& \multicolumn{2}{c}{Any artifact} & \multicolumn{2}{c}{Artifact-free} \\
\cmidrule(lr){2-3}\cmidrule(lr){4-5}
Analysis & $\Delta$Dice [95\% CI] & $p$ & $\Delta$Dice [95\% CI] & $p$ \\
\midrule
Primary (top-1, seed-averaged, canvas) & \res & \res & \res & \res \\
Union rule, confidence $\ge0.25$ & \res & \res & \res & \res \\
Seed 0 only & \res & \res & \res & \res \\
Seed 1 only & \res & \res & \res & \res \\
Seed 2 only & \res & \res & \res & \res \\
Bootstrap seeds 1--5 (range of interval limits) & \res & -- & \res & -- \\
Without the 46 test duplicates of training images & \res & \res & \res & \res \\
Without all 50 flagged test images & \res & \res & \res & \res \\
Labels of annotator 1 only (common subset) & \res & \res & \res & \res \\
Labels of annotator 2 only (common subset) & \res & \res & \res & \res \\
\TBD{$\lambda=0.5$ / $2$ (if run)} & \res & \res & \res & \res \\
\bottomrule
\end{tabular}
\pairedcaveat{A1}
\end{table}

% =================================================================================================
\section{Computational Characteristics}\label{sec:s-cost}
Parameters and GFLOPs are measured on the fused models (convolution and batch normalization merged) for one $640\times640$ input with one class. Latency: \texttt{scripts/benchmark.py}, forward pass only, batch 1, FP16, 30 warm-up and 200 timed runs, on \TBD{GPU, driver, CUDA, PyTorch}.

\begin{table}[H]
\centering
\caption{Cost per model. Parameters in millions; GFLOPs at $640\times640$.}
\label{tab:s-cost}
\small
\begin{tabular}{lcccc}
\toprule
Model & Params (M) & GFLOPs & Latency (ms) & FPS \\
\midrule
YOLO26n-seg (stock) & 2.689 & 8.96 & \res & \res \\
A0: + P2 & 2.735 & 11.20 & \res & \res \\
A1: + P2 + CBAM & 2.738 & 11.21 & \res & \res \\
A2: A1 + artifact head & 2.759 & 12.26 & \res & \res \\
\textbf{A4: ArtiGate-YOLO} & 2.759 & 12.26 & \res & \res \\
A3, A5: as A4 & 2.759 & 12.26 & \res & \res \\
A7: AGA after B2 only & 2.759 & 12.26 & \res & \res \\
\bottomrule
\end{tabular}
\end{table}

\noindent Compared with A1, the artifact head and the two AGA blocks add 0.021~M parameters, or 0.8\% of A4, and 1.05~GFLOPs. Nearly all of it is the head, which runs at stride~4 on a $160\times160$ grid: A2, which adds only the head to A1, already has the cost of A4, and replacing CBAM by AGA adds one input channel to each $7\times7$ convolution and the scalar $\beta$.

% =================================================================================================
\section{Seed-Level Results}\label{sec:s-seeds}
\begin{table}[H]
\centering
\caption{Per-seed results on the ISIC 2018 test set: Dice (mean over images), best epoch, and the learned $\beta_2/\beta_3$. Mean $\pm$ SD of the three per-seed means in the last column.}
\label{tab:s-seeds}
\footnotesize
\setlength{\tabcolsep}{3.5pt}
\begin{tabular}{lcccccccc}
\toprule
& \multicolumn{2}{c}{Seed 0} & \multicolumn{2}{c}{Seed 1} & \multicolumn{2}{c}{Seed 2} & & \\
\cmidrule(lr){2-3}\cmidrule(lr){4-5}\cmidrule(lr){6-7}
Row & Dice & Epoch & Dice & Epoch & Dice & Epoch & $\beta_2/\beta_3$ (s0, s1, s2) & Dice mean $\pm$ SD \\
\midrule
A0 & \res & \res & \res & \res & \res & \res & -- & \res \\
A1 & \res & \res & \res & \res & \res & \res & -- & \res \\
A1s & \res & \res & \res & \res & \res & \res & -- & \res \\
A2 & \res & \res & \res & \res & \res & \res & -- & \res \\
A3 & \res & \res & \res & \res & \res & \res & \res & \res \\
\textbf{A4} & \res & \res & \res & \res & \res & \res & \res & \res \\
A5 & \res & \res & \res & \res & \res & \res & \res & \res \\
A7 & \res & \res & \res & \res & \res & \res & \res & \res \\
\bottomrule
\end{tabular}
\end{table}

% =================================================================================================
\section{Published Results of Related Methods}\label{sec:s-published}
The values below are copied from the cited papers and are \emph{not} comparable with ours: splits, input sizes, preprocessing, mask rules and metric implementations differ. They are listed only to place our numbers. \TBD{Fill only with values printed in the papers, with the table or page they come from.}
\begin{table}[H]
\centering
\caption{Published ISIC 2018 segmentation results of related methods (as reported).}
\label{tab:s-published}
\footnotesize
\begin{tabularx}{\textwidth}{@{}l>{\raggedright\arraybackslash}Xcccc@{}}
\toprule
Method & Split / input as reported & Dice & IoU & Params (M) & Source \\
\midrule
FSSM-ResNet~\cite{raikhankyzy2026fssm} & \TBD{} & \res & \res & \res & \TBD{table} \\
Artifact-aware segmenter~\cite{gong2026artifactaware} & \TBD{} & \res & \res & \res & \TBD{table} \\
Dual-resolution network~\cite{singh2025dualresolution} & \TBD{} & \res & \res & \res & \TBD{table} \\
MALUNet~\cite{ruan2022malunet} & \TBD{} & \res & \res & \res & \TBD{table} \\
EGE-UNet~\cite{ruan2023egeunet} & \TBD{} & \res & \res & \res & \TBD{table} \\
\bottomrule
\end{tabularx}
\end{table}

% =================================================================================================
\section{Implementation Details}\label{sec:s-impl}
\subsection{Software and hardware}
Ultralytics 8.4.13, PyTorch \TBD{version}, CUDA \TBD{version}, Python \TBD{version}, on \TBD{GPU}. Code: \TBD{repository URL and commit}. Every run writes its training settings, seed, package versions and GPU, the git commit, SHA-256 hashes of the model configuration and the pretrained weights, and the weight-transfer report to \texttt{provenance.json}.

\subsection{Training schedule}
\begin{table}[H]
\centering
\caption{Training settings shared by all YOLO models (\texttt{configs/train/default.yaml}); unlisted settings keep the Ultralytics 8.4.13 defaults.}
\label{tab:s-hyp}
\small
\begin{tabular}{ll@{\qquad}ll}
\toprule
Setting & Value & Setting & Value \\
\midrule
Epochs / patience & 200 / 100 & Mosaic (off for last) & 1.0 (10 epochs) \\
Batch / input size & 16 / 640 & Mix-up / copy-paste & 0.1 / 0.1 \\
Optimizer & AdamW, lr 0.002, $\beta_1$ 0.9$^{\dagger}$ & Rotation & $\pm180^\circ$ \\
LR schedule & cosine, 3 warm-up epochs & Scale & $\pm0.5$ \\
Mask ratio & 2 & Flips (up-down / left-right) & 0.5 / 0.5 \\
Deterministic & yes & Seeds & 0, 1, 2 \\
Artifact loss weight $\lambda$ & 1 (A2, A4, A5, A7); 0 (A3) & $p_{\mathrm{synth}}$ & 0.7 (A1s--A4, A7); 0 (A5) \\
Artifact map detached & yes; no in A3 & & \\
\bottomrule
\end{tabular}
\par\smallskip{\footnotesize $^{\dagger}$Fixed for every row, with no bias warm-up. This is what the Ultralytics 8.4.13 automatic rule chooses for one-view ISIC 2018 training (8,200 iterations of the nominal batch of 64); above 10,000 iterations the rule chooses MuSGD with lr 0.01, which the final conference run used and which two-view and HAM10000 training would get.\par}
\end{table}

\subsection{Architecture and weight transfer}
\begin{table}[H]
\centering
\caption{Backbone of A4 at scale n (\texttt{configs/models/yolo26n-seg-p2-artigate.yaml}). Layers 3, 4 and 7 are inserted; all other layers are those of YOLO26n-seg, and the neck adds the P2 level. Channels after width scaling.}
\label{tab:s-arch}
\footnotesize
\setlength{\tabcolsep}{4pt}
\begin{tabularx}{\textwidth}{@{}cl>{\raggedright\arraybackslash}Xl@{}}
\toprule
Layer & From & Module & Output \\
\midrule
0 & input & Conv $3\times3$, stride 2 & 16, stride 2 \\
1 & 0 & Conv $3\times3$, stride 2 & 32, stride 4 \\
2 & 1 & C3k2 (B2) & 64, stride 4 \\
3 & 2 & ArtifactHead ($c_h=32$, $K=4$) & 4 logits, stride 4 \\
4 & 2, 3 & AGA ($7\times7$, full mode, $\bar{A}$ detached except in A3) & 64, stride 4 \\
5 & 4 & Conv $3\times3$, stride 2 & 64, stride 8 \\
6 & 5 & C3k2 (B3) & 128, stride 8 \\
7 & 6, 3 & AGA & 128, stride 8 \\
8--13 & & Conv, C3k2, Conv, C3k2, SPPF, C2PSA & 128--256, strides 16--32 \\
14--31 & & PAN neck with P2 level & P2 32, P3 64, P4 128, P5 256 \\
32 & 22, 25, 28, 31 & Segment26 with P2 input (32 mask coefficients, prototype branch of 64 channels) & boxes P2--P5, prototypes stride 4 \\
\bottomrule
\end{tabularx}
\end{table}

\noindent Channel counts were read from a forward pass of the built model. AGA and the artifact head take their input channels from the model parser: Ultralytics' \texttt{parse\_model} is extended in memory with one branch for these modules (\texttt{artigate.modules.register\_modules}), so the configuration files list only their own arguments. Weights of the COCO-pretrained YOLO26n-seg are transferred through an explicit map from the 17 stock backbone and top-down neck layers (0--16) to the matching layers of each configuration, skipping inserted layers; every mapped pair must agree in module type, repeats and arguments, and tensors are copied only if their shapes match. For A4, 348 tensors are transferred. The layers after the top-down path (the P2 level, the bottom-up path and the segmentation head) are initialized randomly. Resuming a run loads the full checkpoint of the same architecture.

\subsection{Determinism}
The channel-attention max pooling of CBAM and AGA is computed with \texttt{amax} over the spatial dimensions instead of adaptive max pooling, whose CUDA backward pass is not deterministic; the forward outputs of the two agree to within $10^{-6}$ (unit test). Training runs with \texttt{deterministic=True}. Synthetic artifacts draw their seed from NumPy's global random state, which Ultralytics seeds in every dataloader worker, so they are as reproducible as the other augmentations. \TBD{Confirm by repeating one short run with the same seed.}

\subsection{Leakage checks}
Before scoring, each test set is compared with the training data (\texttt{scripts/check\_leakage.py}): shared ISIC identifiers, and near-duplicate candidates from a 64-bit difference hash (Hamming distance $\le4$) confirmed by a correlation $\ge0.95$ of $64\times64$ grey thumbnails of the image content. Rows and columns that consist only of letterbox padding are cropped before the correlation: on whole canvases, 111 unrelated ISIC 2018 / HAM10000 pairs exceeded the threshold only because photos with the same aspect ratio share identical padding bands. On the content, the highest correlation between ISIC 2018 training images and HAM10000 was 0.944, and the top pairs, checked by eye, show different lesions. External test sets are cleaned on the test side; for the official ISIC 2018 validation and test partitions, pairs are reported and a sensitivity analysis without them is given, but the partitions are kept so that results stay comparable with other work. Every candidate pair with a content correlation between 0.90 and 0.99 was checked by eye, because the 0.95 threshold missed true copies (one test image is a training photograph with a watermark); pairs at 0.99 or above are near-identical. Result (\texttt{data/splits/isic2018\_test\_near\_duplicates.csv}): 5 validation images, all the same photograph as a training image; 50 test images, of which 43 show the same photograph, 3 the same lesion photographed again and 4 are uncertain. Because the best checkpoint is chosen on the validation set, its 5 duplicates make that choice slightly optimistic for every model alike. Near-duplicates detect the same photograph, not different photographs of the same lesion; ISIC 2018 does not publish lesion identifiers, so the latter cannot be excluded (main text, limitations).

\begin{table}[H]
\centering
\caption{Leakage check results.}
\label{tab:s-leakage}
\small
\begin{tabular}{lccc}
\toprule
Test set vs.\ ISIC 2018 train & Shared IDs & Likely duplicates & Removed \\
\midrule
ISIC 2018 validation & 0 & 5 & kept \\
ISIC 2018 test & 0 & 46 (+4 uncertain) & kept$^{\dagger}$ \\
HAM10000 (one image per lesion) & 0 & 0 & 0 \\
PH2 & 0 & 0 & 0 \\
DDI (FEDD subset) & 0 & 0 & 0 \\
\bottomrule
\end{tabular}
\par\smallskip{\footnotesize $^{\dagger}$Kept in the primary analysis; left out in the sensitivity analyses of Table~\ref{tab:s-sensitivity}.\par}
\end{table}

\subsection{Non-YOLO baselines}
\TBD{For U-Net, MALUNet, EGE-UNet and nnU-Net 2D: code source and commit, input size, augmentation, epochs, optimizer, and how their outputs are mapped to the $640\times640$ canvas for scoring with the same metric code.}

% =================================================================================================
\section{Dataset Inventory and Split Manifests}\label{sec:s-data}
\begin{table}[H]
\centering
\caption{Data used, with licences and the files that fix each split (released with the code).}
\label{tab:s-data}
\footnotesize
\begin{tabularx}{\textwidth}{@{}l>{\raggedright\arraybackslash}X>{\raggedright\arraybackslash}X>{\raggedright\arraybackslash}X@{}}
\toprule
Dataset & Images used & Licence & Manifest \\
\midrule
ISIC 2018 Task 1~\cite{codella2019isic2018} & 2,594 / 100 / 1,000 (official) & \TBD{CC BY-NC 4.0, verify} & \TBD{data/splits/...} \\
HAM10000~\cite{tschandl2018ham10000} & 7,470, one image per lesion (lowest image identifier): nv 5,403, bkl 727, mel 614, bcc 327, akiec 228, vasc 98, df 73; external test & \TBD{CC BY-NC 4.0, verify} & \TBD{} \\
PH2~\cite{mendonca2013ph2} & 200 (test) & \TBD{terms of use} & \TBD{} \\
Artifact labels, ISIC 2018 train~\cite{bissoto2020debiasing} & 2,594 & \TBD{} & \nolinkurl{data/splits/isic2018_train_flags.csv} \\
Artifact labels, ISIC 2018 test (ours) & 1,000 & \TBD{to be released} & \TBD{\nolinkurl{data/splits/isic2018_test_artifacts.csv}} \\
DDI with FEDD masks~\cite{daneshjou2022ddi,carrion2023fedd} & 333 of 656 (those with FEDD masks); external test & \TBD{Stanford AIMI terms of use; FEDD repository licence} & \TBD{} \\
\bottomrule
\end{tabularx}
\end{table}

% =================================================================================================
\section{Relation to the Conference Version}\label{sec:s-conference}
This work extends our conference paper~\cite{iceeis2026own}, which introduced the YOLO26n-seg model with a P2 level and CBAM (row A1) and compared input views with and without DullRazor preprocessing. The present manuscript adds the artifact head, AGA, the partial-label training targets, the artifact-stratified evaluation, the external test sets and the image-level statistical protocol. During this work we found that the conference code loaded the pretrained weights before Ultralytics rebuilt the model from its configuration, so the conference model was trained from random initialization; all A1 numbers in this manuscript come from new runs with the corrected transfer and with AdamW instead of the conference run's MuSGD (Table~\ref{tab:s-hyp}), and therefore differ from the conference paper. The same holds for the DullRazor view (v7): hair removal now runs at the 640-pixel scale for which its parameters were chosen, whereas the conference implementation ran it at the original resolution, where it removed little hair. \TBD{Text overlap with the conference paper (similarity report) and the journal's policy on extended conference papers.}

\clearpage
\begingroup
\small
\bibliographystyle{IEEEtran}
\bibliography{refs}
\endgroup

\end{document}
